% page 556 du fac-simile (image p-555 du scan)
% bloc 157.5 x 242.0 mm ; marge gauche 8.0 mm
\pgc{-12.700mm}{0.000mm}{
\l{\cel{3}548.}
\l{}
\l{\sou{sporno}. Brancho ek metalo, qua esas muntita an la talono di}
\l{\cel{3}la kavalkanto e finas per disko dentetoza e movebla, quan}
\l{\cel{3}la kavalkanto sinkas aden la flanki di sua kavalo kom}
\l{\cel{3}stimulo. - DFI.}
\l{}
\l{\sou{sporofito}. (\sou{biol.}) Stando di la ovo di la \textquotedbl{}uridinei\textquotedbl{} lor la}
\l{\cel{3}fekundigeso, kande la du pronuklei ne esas ja kunfuzita.}
\l{}
\l{\sou{sportar}.\cel{2}(\sou{netrans.})\cel{2}Exercar su korpale, metodoze, por de-}
\l{\cel{3}velopar en su la bona qualesi di sua korpo ed ula bona}
\l{\cel{3}qualesi di la psiko (loyaleso, energiozeso, perseverem-}
\l{\cel{3}eso) per pedo-balono, natado, skermado, aviacado, e c.}
\l{\cel{3}- DEFIRS.}
\l{}
\l{\sou{spoto}. (\sou{fiziol}.) Mikra saliajo, o mikra makulo sur la pelo,}
\l{\cel{3}qua igas -ulakaze - plu \textquotedbl{}spicoza\textquotedbl{} la fizionomio (di mu-}
\l{\cel{3}liero). - E.}
\l{}
\l{\sou{spozo}. Persono unionita a persono altra-sexua per mariajo.}
\l{\cel{3}EFIS.}
\l{}
\l{\sou{sprato}. (\sou{zool.}) Nomo vulgara di mikra fisho, sorto di haringo,}
\l{\cel{3}qua vivas proxim la litoro Franca di Atlantiko. -}
\l{\cel{3}L.\cel{1}\sou{clupea sprattus}. -}
\l{}
\l{\sou{spricar}. (\sou{netrans.}) I. (\sou{aludante liquido, fluido)}. Springar}
\l{\cel{3}impetuoze, ek orifico relative mikra. - II. (\sou{metaf.})}
\l{\cel{3}Springar, des-intrikante su, ekirar subite, vivace. -D}
\l{}
\l{\sou{sproso}.\cel{2}(\sou{bot.)}\cel{2}Kresko komencala diburjono. - D.}
\l{}
\l{\sou{spular}. (\sou{trans.}) Volvar (kozo) cirkum altra quaze filo sur}
\l{\cel{3}bobino. DE.}
\l{}
\l{\sou{spumo}. I.Quaza moso qua naskas surface di liquido quan onu}
\l{\cel{3}agitas, varmigas o qua fermentacas. - II. Bavajo mosatra}
\l{\cel{3}qua naskas ye la boko di ula animali kande li esas eci-}
\l{\cel{3}tita od iritita. - EFIS.}
\l{}
\l{\sou{sputar}. (\sou{trans}.)De-jetar (to quon onu havas en sua boko).}
\l{\cel{3}EFI.}
\l{}
\l{\sou{squadro}. Instrumento ek ligno, metalo, quan onu uzas por tra-}
\l{\cel{3}sar ortanguli o perpendikli. - IS.}
\l{}
\l{\sou{squalo}. (\sou{zool.)}\cel{2}Genero de fishi kartilagoza. - L.\cel{1}\sou{squalus}.}
\l{\cel{3}FIS.}
\l{}
\l{\sou{squamo}.\cel{1}\sou{(zool}.) Singla ek la plaki inter-apudpozita od imbri-}
\l{\cel{3}kita qui tegas la pelo di ula fishi, di ula repteri, e c.}
\l{}
\l{\sou{squatar}. (\sou{netrans.}) I. (\sou{aludante la animali}).Sidar sur sua}
\l{\cel{3}gropo. - II. (\sou{aludante la homo)}. Sidar tale ke la gambi}
\l{\cel{3}esas rifaldita, la glutei tushante la taloni. - E.}
}
